\documentclass[10pt,a4paper]{article}
\usepackage[margin=0.8in]{geometry}
\usepackage{amsmath,amssymb,hyperref}
\hypersetup{colorlinks=true,urlcolor=blue,linkcolor=blue}
\setlength{\parindent}{0pt}\setlength{\parskip}{6pt}
\begin{document}
\begin{center}
{\Large Infinitely many integral figure-eight characters}\\[4pt]
TLMC 00000001734 \quad lizaixi01 \quad 24 October 2026
\end{center}

\textbf{The claim and its scope.}
The source asserts finiteness of the integral points of the $\mathrm{SL}_2$
character variety of a fibered hyperbolic $3$-manifold group. It imposes no
irreducibility, canonical-component, or fixed-boundary-trace restriction.
We disprove this whole-variety assertion using the figure-eight knot exterior
and an infinite family of actual integral trace characters. We make no claim
against versions with any of those additional restrictions, or about the
source's undefined Hitchin trace-field phrase.

\textbf{Published geometric and character identifications.}
We use the standard figure-eight identification in \cite{hmp}, printed p.~2:
its exterior fibers over a circle with once-punctured-torus fiber and has group
\[
 \Gamma=\langle a,b,t\mid tat^{-1}=ab,\quad tbt^{-1}=bab\rangle.
 \tag{1}
\]
Its complete hyperbolic structure is also discussed on printed p.~3.
A character is the full function $\chi_\rho(g)=\operatorname{tr}(\rho(g))$
of a representation $\rho:\Gamma\to\mathrm{SL}_2(\mathbb C)$; see also
\cite{porti}, Definition~1.3 and Corollary~1.6, printed pp.~II--2--II--3.
These published identifications of the geometric group and character variety
are inputs to interpreting the formal example. The Lean development does
\emph{not} reprove the topology, the hyperbolic structure, or general geometric
invariant theory. It constructs the precise presented group (1) and proves
the representation, character, integrality and infinitude claims below.

\textbf{Genuine integral representations.}
For $n\in\mathbb N$, put $m_n=n+3$ and define
\[
 A_n=\begin{pmatrix}m_n&-1\\1&0\end{pmatrix},\qquad
 \rho_n(a)=\rho_n(b)=I_2,\qquad \rho_n(t)=A_n.
 \tag{2}
\]
Since $\det A_n=1$, these matrices lie in $\mathrm{SL}_2(\mathbb Z)$.
Both relators in (1) map to the identity, so the universal property of the
presented group yields an actual group homomorphism
$\rho_n:\Gamma\to\mathrm{SL}_2(\mathbb Z)$.
The native scalar-extension homomorphism gives a complex representation
$\rho_n^{\mathbb C}$. The construction is not just a list of proposed
coordinate values: it defines the image of every group element.

\textbf{All traces are integral, and the characters are distinct.}
For every $g\in\Gamma$, the proof establishes
\[
 \chi_n(g)=\operatorname{tr}(\rho_n^{\mathbb C}(g))
    =\iota\bigl(\operatorname{tr}(\rho_n(g))\bigr)\in\iota(\mathbb Z),
 \qquad \chi_n(t)=n+3, \tag{3}
\]
where $\iota:\mathbb Z\hookrightarrow\mathbb C$ is the usual embedding.
Thus the full character is integral on every word, not only on the three
generators. Evaluation at $t$ shows that $n\mapsto\chi_n$ is injective.
Consequently the set of integral characters among all actual complex
representations of $\Gamma$ is infinite. Trace invariance under conjugation
is also proved, so the different characters cannot arise from conjugate
representations.

\newpage
\textbf{Connection to the published integral coordinate model.}
Proposition~3.1 of \cite{hmp}, printed p.~7, describes the full
$\mathrm{SL}_2$ character variety in coordinates
\[
 (x,y,z)=(\chi(a),\chi(b),\chi(t)).
\]
It is the union of the line $x=y=2$ and the component
$(x-1)(y-1)=1$, $z^2=x+y+1$. In particular the abelian line belongs to
the \emph{full} variety. The Lean proof works with the actual integer
polynomial zero locus
\[
 \begin{split}
 F_0&=(X-1)(Y-1)-1,\\
 F_1&=(X-2)(Z^2-X-Y-1),\\
 F_2&=(Y-2)(Z^2-X-Y-1)
 \end{split} \tag{4}
\]
and proves, over $\mathbb Z$, the precise set-theoretic equivalence
\[
 \bigl(F_0=F_1=F_2=0\bigr)\iff
 \bigl(x=y=2\bigr)\ \lor\
 \bigl((x-1)(y-1)=1\ \land\ z^2=x+y+1\bigr).
\]
No claim of a scheme-theoretic ideal computation is needed.
The genuine trace coordinates of (2) are
\[
 P_n=(2,2,n+3)\in\mathbb Z^3. \tag{5}
\]
The formal proof verifies $F_i(P_n)=0$, proves $n\mapsto P_n$ injective,
and supplies the representation realizing every $P_n$, together with (3).
It therefore proves infinitude even for the subset of integer points in (4)
realized by actual representations whose \emph{every} word has integral trace.
This links the infinite family to the published character-variety coordinates;
an arbitrary integer line has not been labeled a character variety.

\textbf{Consequences and limits.}
The figure-eight exterior has one torus boundary component. Infinitely many
distinct integral points already contradict the source's necessary finiteness
clause, so no particular formula for the proposed boundary-count bound is
needed. The family lies on the abelian component and does not establish
infinitude on the irreducible or canonical component. It also varies the
$t$-trace, so it is not a counterexample to a fixed-boundary-character variant.
These restrictions were not present in the printed finiteness claim.

\textbf{Formal coverage and verification.}
\texttt{Main.lean} uses the native free and presented groups,
$\mathrm{SL}_2$ subgroups and scalar extension, matrix trace,
multivariate polynomials and evaluation, and actual sets of full trace
functions. Its final theorem proves infinitude of realized integral points.
The project pins Lean~4.33.0 and Mathlib
\texttt{db584cd6d46c92f209a44c0f1c829460d327499d}.
The accompanying evidence records the source-bound build, ordinary
\texttt{\#print axioms} for the complete namespace, and the official
fresh-kernel replay. There are no admitted proofs, additional axioms,
or uses of \texttt{native\_decide}. Production timings are not controlled
research-efficiency evidence.

\begin{thebibliography}{9}
\bibitem{hmp} M. Heusener, V. Mu\~noz and J. Porti,
\emph{The $\mathrm{SL}(3,\mathbb C)$-character variety of the figure eight knot},
\href{https://arxiv.org/abs/1505.04451}{arXiv:1505.04451}, v1 (2015).
Printed pp.~2--3 and Proposition~3.1 on p.~7.
\bibitem{porti} J. Porti, \emph{Character varieties and knot symmetries},
Winter Braids Lecture Notes \textbf{4} (2017), Course II, 1--21.
\href{https://doi.org/10.5802/wbln.18}{doi:10.5802/wbln.18}.
Definition~1.3, Corollary~1.6 and Examples~2.6, 2.9.
\end{thebibliography}
\end{document}
